\section{BFS multiorigen}
\whenuse{Distancia al origen m\'as cercano o expansi\'on simult\'anea desde muchas
fuentes en un grafo sin pesos.}

Usar cuando se pide la distancia al origen m\'as cercano, expansi\'on simult\'anea,
Voronoi sobre un grafo no pesado o el instante en que llega primero alguna fuente.
Es una BFS normal inicializando la cola con todas las fuentes a distancia cero.
Los comentarios opcionales agregan padres y \texttt{owner[v]}, la fuente
m\'as cercana. En empates gana la primera de \texttt{sources}; los inalcanzables
tienen distancia y fuente $-1$. Sin fuentes, todos quedan inalcanzables.

\complexity{tiempo $O(n+m)$; memoria $O(n)$.}

\begin{lstlisting}[style=cpp]
vector<int> multi_source_bfs(const vector<vector<int>> &graph,
                             const vector<int> &sources
    // , vector<int>& parent, vector<int>& owner // Reconstruccion y fuente.
) {
    int node_count = (int)graph.size();
    vector<int> distance(node_count, -1);
    // parent.assign(node_count, -1); owner.assign(node_count, -1);
    queue<int> pending;
    for (int source : sources) if (distance[source] == -1) {
        distance[source] = 0;
        // owner[source] = source; // Cada fuente se pertenece a si misma.
        pending.push(source);
    }
    while (!pending.empty()) {
        int vertex = pending.front(); pending.pop();
        for (int to : graph[vertex]) if (distance[to] == -1) {
            distance[to] = distance[vertex] + 1;
            // parent[to] = vertex; owner[to] = owner[vertex];
            pending.push(to);
        }
    }
    return distance;
}
\end{lstlisting}

Activar las cuatro l\'ineas comentadas agrega $O(V)$ memoria. Recuperar un
camino hasta la fuente ganadora cuesta $O(V)$:
\begin{lstlisting}[style=cpp]
// void multi_source_path_example() {
//     vector<vector<int>> graph = {{1}, {0,2}, {1,3}, {2,4}, {3}};
//     vector<int> parent, owner;
//     auto distance = multi_source_bfs(graph, {0,4}, parent, owner);
//     int target = 2;
//     cout << owner[target] << '\n'; // 0: empate ganado por la primera fuente.
//     vector<int> path;
//     if (distance[target] != -1) {
//         for (int v = target; v != -1; v = parent[v]) path.push_back(v);
//         reverse(path.begin(), path.end());
//     }
//     for (int v : path) cout << v << ' '; // 0 1 2
// }
\end{lstlisting}
